% MỞ ĐẦU (LỜI MỞ ĐẦU - Chương 0)

\chapter*{MỞ ĐẦU} % Tên của chương
\addcontentsline{toc}{chapter}{MỞ ĐẦU} % Thêm tên chương vào mục lục

\label{Chapter0} % Để trích dẫn chương này ở chỗ nào đó trong bài, hãy sử dụng lệnh \ref{Chapter0} 

%----------------------------------------------------------------------------------------

Trong bối cảnh công nghệ trí tuệ nhân tạo (AI) và thị giác máy tính đang phát triển
mạnh mẽ, các ứng dụng hỗ trợ con người trong đời sống ngày càng được quan tâm và
nghiên cứu rộng rãi. Đặc biệt, đối với người khiếm thị, việc di chuyển và nhận biết môi
trường xung quanh luôn là một thách thức lớn do hạn chế về khả năng quan sát. Những
khó khăn trong việc phát hiện vật cản, xác định khoảng cách hay nhận biết nguy hiểm
có thể ảnh hưởng trực tiếp đến sự an toàn và khả năng sinh hoạt độc lập của người
khiếm thị trong cuộc sống hằng ngày.

Hiện nay, nhiều giải pháp hỗ trợ người khiếm thị đã được nghiên cứu và phát triển,
từ các thiết bị cảm biến siêu âm truyền thống cho đến các hệ thống ứng dụng trí tuệ nhân
tạo và học sâu (Deep Learning). Trong đó, các phương pháp sử dụng thị giác máy tính
kết hợp với nhận diện đối tượng và ước lượng khoảng cách đang cho thấy nhiều tiềm
năng nhờ khả năng xử lý thông tin trực quan nhanh chóng, chính xác và có thể hoạt động
trong thời gian thực. Tuy nhiên, việc xây dựng các phương pháp có độ chính xác cao,
khả năng thích nghi tốt với môi trường thực tế và đáp ứng yêu cầu xử lý thời gian thực
vẫn còn là một bài toán cần được nghiên cứu và cải thiện.

Xuất phát từ thực tế đó, nhóm em lựa chọn đề tài “Phương pháp nhận diện vật cản và ước
lượng khoảng cách để hỗ trợ người khiếm thị”. Đề tài tập trung nghiên cứu các phương
pháp nhận diện vật cản dựa trên mô hình học sâu kết hợp với kỹ thuật ước lượng khoảng
cách nhằm hỗ trợ người khiếm thị trong quá trình di chuyển. Thông qua việc phân tích,
xây dựng và đánh giá các mô hình phù hợp, đề tài hướng tới việc nâng cao khả năng
nhận biết môi trường xung quanh, hỗ trợ cảnh báo vật cản và tăng độ an toàn cho người
khiếm thị.

Trong báo cáo này, nhóm em trình bày quá trình nghiên cứu lý thuyết, xây dựng phương
pháp, triển khai mô hình và đánh giá kết quả thực nghiệm cho bài toán nhận diện vật cản
và ước lượng khoảng cách nhằm hỗ trợ người khiếm thị. Báo cáo tập trung phân tích
cách thức hoạt động của mô hình, khả năng nhận diện vật cản cũng như hiệu quả ước
lượng khoảng cách trong các điều kiện thực tế khác nhau. Qua đó, đề tài hướng tới việc
đề xuất một phương pháp hỗ trợ người khiếm thị trong quá trình di chuyển, góp phần
nâng cao tính an toàn và khả năng ứng dụng của các giải pháp hỗ trợ thông minh trong
thực tế.